% ══════════════════════════════════════════════════════════════
\section{CENA 5 --- Conectando os Dois Mundos \hfill \normalfont\small\color{corTempo}18:00 -- 19:30}
% ══════════════════════════════════════════════════════════════

\tempo{18:00 -- 18:45}

\begin{visual}
\faFilm\ Split-screen final: heatmap clássico à esquerda, Esfera de Bloch à direita.
Seta conectando $T_{\text{média}} = 66{,}28$ {\textdegree}C ao raio $r_T = 0{,}2752$.
\end{visual}

\begin{fala}
Olha o que construímos. Pegamos uma função de temperatura, integramos sobre
a superfície do chip -- integral \textbf{dupla} -- e chegamos a 66,28 graus.
Esse número alimentou o modelo quântico, que virou o raio da nuvem.
Integramos sobre o volume da esfera -- integral \textbf{tripla} --
e chegamos a 2,08\% de erro.
\end{fala}

\begin{fala}
Resultado importante: $P_{\text{erro}} = r_T^3$, e $r_T$ depende linearmente
de $T_{\text{média}}$. O erro quântico cresce com o \textbf{cubo} da temperatura.
Pequenos aumentos produzem impactos desproporcionalmente grandes na fidelidade.
\end{fala}

\tempo{18:45 -- 19:30}

\begin{center}
\begin{tikzpicture}
\begin{axis}[
  width=9.5cm, height=5cm,
  xlabel={$T_{\text{média}}$ ({\textdegree}C)},
  ylabel={$P_{\text{erro}}$ (\%)},
  xmin=25, xmax=100, ymin=0, ymax=13,
  grid=major, grid style={gray!25},
  title={Probabilidade de erro $\times$ temperatura média},
  xtick={25,40,56,66,80,100},
  xticklabel style={font=\small},
  yticklabel style={font=\small},
]
\addplot[NavyBlue, thick, domain=25:100, samples=30]
  { ((0.5*(x-25)/75)^3)*100 };
\addplot[only marks, mark=*, mark size=3.5pt, red]
  coordinates {(66.28, 2.08)};
\node[font=\scriptsize, red] at (axis cs:74, 3.3)
  {$66{,}28\,\text{\textdegree C} \to 2{,}08\%$};
\addplot[dashed, red!60, domain=25:100] {1.0}
  node[pos=0.04, above, font=\scriptsize, red!70]{tol.\ $1\%$};
\end{axis}
\end{tikzpicture}
\end{center}

\begin{fala}
Uma redução de apenas dez graus na temperatura média do chip --
de 66 para 56 graus -- já seria suficiente para cruzar o limiar de tolerância.
\textbf{Dez graus. É o que separa um qubit funcional de um que erra demais.}
\end{fala}